\documentclass[10pt,a4paper]{amsart}
\usepackage{amsmath,amssymb,amsthm,graphicx}
\usepackage{verbatim}
\newenvironment{reference}{\comment}{\endcomment}
%\newenvironment{reference}{}{}
\newenvironment{slogan}{\comment}{\endcomment}
\newenvironment{history}{\comment}{\endcomment}
% For commutative diagrams we use Xy-pic
\usepackage[all]{xy}

% We use 2cell for 2-commutative diagrams.
\xyoption{2cell}
% Theorem environments.
%
\theoremstyle{plain}
\newtheorem{theorem}[subsection]{Theorem}
\newtheorem{proposition}[subsection]{Proposition}
\newtheorem{lemma}[subsection]{Lemma}

\theoremstyle{definition}
\newtheorem{definition}[subsection]{Definition}
\newtheorem{example}[subsection]{Example}
\newtheorem{exercise}[subsection]{Exercise}
\newtheorem{situation}[subsection]{Situation}

\theoremstyle{remark}
\newtheorem{remark}[subsection]{Remark}
\newtheorem{remarks}[subsection]{Remarks}

\numberwithin{equation}{subsection}

% Macros
%
\def\lim{\mathop{\mathrm{lim}}\nolimits}
\def\colim{\mathop{\mathrm{colim}}\nolimits}
\def\Spec{\mathop{\mathrm{Spec}}}
\def\Hom{\mathop{\mathrm{Hom}}\nolimits}
\def\Ext{\mathop{\mathrm{Ext}}\nolimits}
\def\SheafHom{\mathop{\mathcal{H}\!\mathit{om}}\nolimits}
\def\SheafExt{\mathop{\mathcal{E}\!\mathit{xt}}\nolimits}
\def\Sch{\mathit{Sch}}
\def\Mor{\mathop{\mathrm{Mor}}\nolimits}
\def\Ob{\mathop{\mathrm{Ob}}\nolimits}
\def\Sh{\mathop{\mathit{Sh}}\nolimits}
\def\NL{\mathop{N\!L}\nolimits}
\def\CH{\mathop{\mathrm{CH}}\nolimits}
\def\proetale{{pro\text{-}\acute{e}tale}}
\def\etale{{\acute{e}tale}}
\def\QCoh{\mathit{QCoh}}
\def\Ker{\mathop{\mathrm{Ker}}}
\def\Im{\mathop{\mathrm{Im}}}
\def\Coker{\mathop{\mathrm{Coker}}}
\def\Coim{\mathop{\mathrm{Coim}}}

% Boxtimes
%
\DeclareMathSymbol{\boxtimes}{\mathbin}{AMSa}{"02}

%
% Macros for moduli stacks/spaces
%
\def\QCohstack{\mathcal{QC}\!\mathit{oh}}
\def\Cohstack{\mathcal{C}\!\mathit{oh}}
\def\Spacesstack{\mathcal{S}\!\mathit{paces}}
\def\Quotfunctor{\mathrm{Quot}}
\def\Hilbfunctor{\mathrm{Hilb}}
\def\Curvesstack{\mathcal{C}\!\mathit{urves}}
\def\Polarizedstack{\mathcal{P}\!\mathit{olarized}}
\def\Complexesstack{\mathcal{C}\!\mathit{omplexes}}
% \Pic is the operator that assigns to X its picard group, usage \Pic(X)
% \Picardstack_{X/B} denotes the Picard stack of X over B
% \Picardfunctor_{X/B} denotes the Picard functor of X over B
\def\Pic{\mathop{\mathrm{Pic}}\nolimits}
\def\Picardstack{\mathcal{P}\!\mathit{ic}}
\def\Picardfunctor{\mathrm{Pic}}
\def\Deformationcategory{\mathcal{D}\!\mathit{ef}}
\usepackage{cleveref}
\crefname{originalreference}{original reference}{original references}
\Crefname{originalreference}{Original reference}{Original references}
\makeatletter
\newlabel{algebraic-lemma-map-space-into-stack}{{algebraic-lemma-map-space-into-stack}{source}{}{sourcepointer-1}{}}
\newlabel{algebraic-lemma-map-space-into-stack@cref}{{[originalreference][0][]algebraic-lemma-map-space-into-stack}{[1][0][]source}{}{}{}}
\makeatother
\crefname{originalreference}{original reference}{original references}
\Crefname{originalreference}{Original reference}{Original references}
\makeatletter
\newlabel{algebraic-remark-flat-fp-presentation}{{algebraic-remark-flat-fp-presentation}{source}{}{sourcepointer-2}{}}
\newlabel{algebraic-remark-flat-fp-presentation@cref}{{[originalreference][0][]algebraic-remark-flat-fp-presentation}{[1][0][]source}{}{}{}}
\makeatother
\crefname{originalreference}{original reference}{original references}
\Crefname{originalreference}{Original reference}{Original references}
\makeatletter
\newlabel{criteria-theorem-flat-groupoid-gives-algebraic-stack}{{criteria-theorem-flat-groupoid-gives-algebraic-stack}{source}{}{sourcepointer-3}{}}
\newlabel{criteria-theorem-flat-groupoid-gives-algebraic-stack@cref}{{[originalreference][0][]criteria-theorem-flat-groupoid-gives-algebraic-stack}{[1][0][]source}{}{}{}}
\makeatother
\crefname{originalreference}{original reference}{original references}
\Crefname{originalreference}{Original reference}{Original references}
\makeatletter
\newlabel{spaces-groupoids-lemma-quotient-stack-restrict}{{spaces-groupoids-lemma-quotient-stack-restrict}{source}{}{sourcepointer-4}{}}
\newlabel{spaces-groupoids-lemma-quotient-stack-restrict@cref}{{[originalreference][0][]spaces-groupoids-lemma-quotient-stack-restrict}{[1][0][]source}{}{}{}}
\makeatother
\crefname{originalreference}{original reference}{original references}
\Crefname{originalreference}{Original reference}{Original references}
\makeatletter
\newlabel{spaces-morphisms-lemma-noetherian-finite-type-finite-presentation}{{spaces-morphisms-lemma-noetherian-finite-type-finite-presentation}{source}{}{sourcepointer-5}{}}
\newlabel{spaces-morphisms-lemma-noetherian-finite-type-finite-presentation@cref}{{[originalreference][0][]spaces-morphisms-lemma-noetherian-finite-type-finite-presentation}{[1][0][]source}{}{}{}}
\makeatother
\crefname{originalreference}{original reference}{original references}
\Crefname{originalreference}{Original reference}{Original references}
\makeatletter
\newlabel{spaces-over-fields-lemma-mono-towards-locally-quasi-finite-over-field}{{spaces-over-fields-lemma-mono-towards-locally-quasi-finite-over-field}{source}{}{sourcepointer-6}{}}
\newlabel{spaces-over-fields-lemma-mono-towards-locally-quasi-finite-over-field@cref}{{[originalreference][0][]spaces-over-fields-lemma-mono-towards-locally-quasi-finite-over-field}{[1][0][]source}{}{}{}}
\makeatother
\crefname{originalreference}{original reference}{original references}
\Crefname{originalreference}{Original reference}{Original references}
\makeatletter
\newlabel{stacks-properties-lemma-monomorphism}{{stacks-properties-lemma-monomorphism}{source}{}{sourcepointer-7}{}}
\newlabel{stacks-properties-lemma-monomorphism@cref}{{[originalreference][0][]stacks-properties-lemma-monomorphism}{[1][0][]source}{}{}{}}
\makeatother
\crefname{originalreference}{original reference}{original references}
\Crefname{originalreference}{Original reference}{Original references}
\makeatletter
\newlabel{stacks-properties-lemma-monomorphism-injective-points}{{stacks-properties-lemma-monomorphism-injective-points}{source}{}{sourcepointer-8}{}}
\newlabel{stacks-properties-lemma-monomorphism-injective-points@cref}{{[originalreference][0][]stacks-properties-lemma-monomorphism-injective-points}{[1][0][]source}{}{}{}}
\makeatother
\crefname{originalreference}{original reference}{original references}
\Crefname{originalreference}{Original reference}{Original references}
\makeatletter
\newlabel{stacks-properties-lemma-points-cartesian}{{stacks-properties-lemma-points-cartesian}{source}{}{sourcepointer-9}{}}
\newlabel{stacks-properties-lemma-points-cartesian@cref}{{[originalreference][0][]stacks-properties-lemma-points-cartesian}{[1][0][]source}{}{}{}}
\makeatother
\crefname{originalreference}{original reference}{original references}
\Crefname{originalreference}{Original reference}{Original references}
\makeatletter
\newlabel{stacks-properties-lemma-residual-gerbe-unique}{{stacks-properties-lemma-residual-gerbe-unique}{source}{}{sourcepointer-10}{}}
\newlabel{stacks-properties-lemma-residual-gerbe-unique@cref}{{[originalreference][0][]stacks-properties-lemma-residual-gerbe-unique}{[1][0][]source}{}{}{}}
\makeatother
\crefname{originaltheorem}{original theorem}{original theorems}
\Crefname{originaltheorem}{Original theorem}{Original theorems}
\makeatletter
\newlabel{theorem-quasi-DM}{{theorem-quasi-DM}{source}{}{sourcepointer-11}{}}
\newlabel{theorem-quasi-DM@cref}{{[originaltheorem][0][]theorem-quasi-DM}{[1][0][]source}{}{}{}}
\makeatother
\title{Every point of a quasi-DM algebraic stack has a residual gerbe}
\author{The Stacks Project Authors}
\date{stacks-project:stacks-morphisms.tex:local-snapshot-20261012; GNU FDL 1.2 or later}
\begin{document}
\maketitle
\noindent\textit{Editorial note.} All delivered mathematical text is the author's own, copied from the pinned Stacks Project chapter source sources/stacks-algebraic-stacks-morphisms/pinned-original-source.tex; no statement or proof step was rewritten, reordered or removed.\par\medskip
\noindent\textit{Editorial note.} Preamble provenance: the Stacks Project keeps one single shared preamble for every chapter; this card takes it directly from the staged sources/stacks-algebraic-stacks-morphisms/preamble.tex as three verbatim spans (lines 8-12, 14-18 and 154-217), so no derived file was written and nothing was added or redefined. The wrapper class and the base packages are supplied by the assembler.\par\medskip
\noindent\textit{Editorial note.} Bibliography provenance: neither the counted statement nor its proof contains a citation (this lane rescans the whole span before writing the card), so the card carries no bibliography and no entry was reconstructed.\par\medskip
\noindent\textit{Editorial note.} Selection provenance: this card was opened under the file-based control contract (CONTROL.md), with the mechanical pre-card gate run on its statement and proof interval before the card was written.\par\medskip
\noindent\textit{Allowed standard prerequisites (editorial).}\begin{itemize}
\item The presentation of a quasi-DM algebraic stack by a flat, locally finitely presented and locally quasi-finite groupoid in algebraic spaces. Invoked at the start of the counted proof (source lines 7095-7097) as a pointer to the author theorem theorem-quasi-DM of this chapter. Provides the groupoid from which the residual gerbe is built.
\item The quotient stack of a flat groupoid, the behaviour of that quotient under restricting the groupoid along a monomorphism, and the injectivity of the induced map on points for a monomorphism of stacks. Invoked inside the counted proof (source lines 7101-7105, 7132-7144) as cross-chapter pointers. Identifies the stack constructed in the proof with the original one and shows that its map to the stack is injective on points.
\item The uniqueness of the residual gerbe, and the criteria for a monomorphism with singleton image and locally Noetherian reduced source to be it. Invoked at the end of the counted proof (source lines 7156-7161) as a pointer to a result of the chapter on properties of stacks. Concludes that the constructed stack is the residual gerbe.
\end{itemize}\medskip
\ifdefined\hypertarget\hypertarget{sourcepointer-1}{}\fi
\noindent\textit{Source reference.} Original reference algebraic-lemma-map-space-into-stack: sources/stacks-algebraic-stacks-morphisms/pinned-original-source.tex, source lines 3831--3831. Navigation pointer only: the counted text refers to the Stacks item algebraic-lemma-map-space-into-stack, which belongs to another chapter of the project and is not part of this card. The pointer carries the target's own label as its number because this chapter's staged snapshot does not contain the other chapter's tag table; its content is not reproduced here.\par\medskip
\ifdefined\hypertarget\hypertarget{sourcepointer-2}{}\fi
\noindent\textit{Source reference.} Original reference algebraic-remark-flat-fp-presentation: sources/stacks-algebraic-stacks-morphisms/pinned-original-source.tex, source lines 4422--4422. Navigation pointer only: the counted text refers to the Stacks item algebraic-remark-flat-fp-presentation, which belongs to another chapter of the project and is not part of this card. The pointer carries the target's own label as its number because this chapter's staged snapshot does not contain the other chapter's tag table; its content is not reproduced here.\par\medskip
\ifdefined\hypertarget\hypertarget{sourcepointer-3}{}\fi
\noindent\textit{Source reference.} Original reference criteria-theorem-flat-groupoid-gives-algebraic-stack: sources/stacks-algebraic-stacks-morphisms/pinned-original-source.tex, source lines 3850--3850. Navigation pointer only: the counted text refers to the Stacks item criteria-theorem-flat-groupoid-gives-algebraic-stack, which belongs to another chapter of the project and is not part of this card. The pointer carries the target's own label as its number because this chapter's staged snapshot does not contain the other chapter's tag table; its content is not reproduced here.\par\medskip
\ifdefined\hypertarget\hypertarget{sourcepointer-4}{}\fi
\noindent\textit{Source reference.} Original reference spaces-groupoids-lemma-quotient-stack-restrict: sources/stacks-algebraic-stacks-morphisms/pinned-original-source.tex, source lines 7138--7138. Navigation pointer only: the counted text refers to the Stacks item spaces-groupoids-lemma-quotient-stack-restrict, which belongs to another chapter of the project and is not part of this card. The pointer carries the target's own label as its number because this chapter's staged snapshot does not contain the other chapter's tag table; its content is not reproduced here.\par\medskip
\ifdefined\hypertarget\hypertarget{sourcepointer-5}{}\fi
\noindent\textit{Source reference.} Original reference spaces-morphisms-lemma-noetherian-finite-type-finite-presentation: sources/stacks-algebraic-stacks-morphisms/pinned-original-source.tex, source lines 3847--3847. Navigation pointer only: the counted text refers to the Stacks item spaces-morphisms-lemma-noetherian-finite-type-finite-presentation, which belongs to another chapter of the project and is not part of this card. The pointer carries the target's own label as its number because this chapter's staged snapshot does not contain the other chapter's tag table; its content is not reproduced here.\par\medskip
\ifdefined\hypertarget\hypertarget{sourcepointer-6}{}\fi
\noindent\textit{Source reference.} Original reference spaces-over-fields-lemma-mono-towards-locally-quasi-finite-over-field: sources/stacks-algebraic-stacks-morphisms/pinned-original-source.tex, source lines 7123--7123. Navigation pointer only: the counted text refers to the Stacks item spaces-over-fields-lemma-mono-towards-locally-quasi-finite-over-field, which belongs to another chapter of the project and is not part of this card. The pointer carries the target's own label as its number because this chapter's staged snapshot does not contain the other chapter's tag table; its content is not reproduced here.\par\medskip
\ifdefined\hypertarget\hypertarget{sourcepointer-7}{}\fi
\noindent\textit{Source reference.} Original reference stacks-properties-lemma-monomorphism: sources/stacks-algebraic-stacks-morphisms/pinned-original-source.tex, source lines 1167--1167. Navigation pointer only: the counted text refers to the Stacks item stacks-properties-lemma-monomorphism, which belongs to another chapter of the project and is not part of this card. The pointer carries the target's own label as its number because this chapter's staged snapshot does not contain the other chapter's tag table; its content is not reproduced here.\par\medskip
\ifdefined\hypertarget\hypertarget{sourcepointer-8}{}\fi
\noindent\textit{Source reference.} Original reference stacks-properties-lemma-monomorphism-injective-points: sources/stacks-algebraic-stacks-morphisms/pinned-original-source.tex, source lines 7144--7144. Navigation pointer only: the counted text refers to the Stacks item stacks-properties-lemma-monomorphism-injective-points, which belongs to another chapter of the project and is not part of this card. The pointer carries the target's own label as its number because this chapter's staged snapshot does not contain the other chapter's tag table; its content is not reproduced here.\par\medskip
\ifdefined\hypertarget\hypertarget{sourcepointer-9}{}\fi
\noindent\textit{Source reference.} Original reference stacks-properties-lemma-points-cartesian: sources/stacks-algebraic-stacks-morphisms/pinned-original-source.tex, source lines 2126--2126. Navigation pointer only: the counted text refers to the Stacks item stacks-properties-lemma-points-cartesian, which belongs to another chapter of the project and is not part of this card. The pointer carries the target's own label as its number because this chapter's staged snapshot does not contain the other chapter's tag table; its content is not reproduced here.\par\medskip
\ifdefined\hypertarget\hypertarget{sourcepointer-10}{}\fi
\noindent\textit{Source reference.} Original reference stacks-properties-lemma-residual-gerbe-unique: sources/stacks-algebraic-stacks-morphisms/pinned-original-source.tex, source lines 3867--3867. Navigation pointer only: the counted text refers to the Stacks item stacks-properties-lemma-residual-gerbe-unique, which belongs to another chapter of the project and is not part of this card. The pointer carries the target's own label as its number because this chapter's staged snapshot does not contain the other chapter's tag table; its content is not reproduced here.\par\medskip
\ifdefined\hypertarget\hypertarget{sourcepointer-11}{}\fi
\noindent\textit{Source reference.} Original theorem theorem-quasi-DM: sources/stacks-algebraic-stacks-morphisms/pinned-original-source.tex, source lines 619--619. Navigation pointer only: the counted text refers to the Stacks item theorem-quasi-DM, which belongs to another chapter of the project and is not part of this card. The pointer carries the target's own label as its number because this chapter's staged snapshot does not contain the other chapter's tag table; its content is not reproduced here.\par\medskip
\begin{lemma}
\label{lemma-every-point-residual-gerbe-quasi-DM}
Let $\mathcal{X}$ be a quasi-DM algebraic stack.
Then the residual gerbe of $\mathcal{X}$ at $x$ exists for
every $x \in |\mathcal{X}|$.
\end{lemma}

\begin{proof}
Choose a scheme $U$ and a surjective, flat, locally finite presented,
and locally quasi-finite morphism $U \to \mathcal{X}$, see
Theorem \ref{theorem-quasi-DM}.
Set $R = U \times_\mathcal{X} U$. The projections $s, t : R \to U$
are surjective, flat, locally of finite presentation, and
locally quasi-finite as base changes of the morphism $U \to \mathcal{X}$.
There is a canonical morphism $[U/R] \to \mathcal{X}$ (see
Algebraic Stacks, Lemma \ref{algebraic-lemma-map-space-into-stack})
which is an equivalence because $U \to \mathcal{X}$ is surjective, flat,
and locally of finite presentation, see
Algebraic Stacks, Remark \ref{algebraic-remark-flat-fp-presentation}.
Thus we may assume that $\mathcal{X} = [U/R]$ where
$(U, R, s, t, c)$ is a groupoid in algebraic spaces such that
$s, t : R \to U$ are surjective, flat, locally of finite presentation,
and locally quasi-finite. Set
$$
U' = \coprod\nolimits_{u \in U\text{ lying over }x} \Spec(\kappa(u)).
$$
The canonical morphism $U' \to U$ is a monomorphism. Let
$$
R' = U' \times_\mathcal{X} U' =
R \times_{(U \times U)} (U' \times U')
$$
Because $U' \to U$ is a monomorphism we see that both projections
$s', t' : R' \to U'$ factor as a monomorphism followed by a locally
quasi-finite morphism. Hence, as $U'$ is a disjoint union of spectra
of fields, using
Spaces over Fields, Lemma
\ref{spaces-over-fields-lemma-mono-towards-locally-quasi-finite-over-field}
we conclude that the morphisms $s', t' : R' \to U'$ are locally quasi-finite.
Again since $U'$ is a disjoint union of spectra of fields, the morphisms
$s', t'$ are also flat. Finally, $s', t'$ locally quasi-finite
implies $s', t'$ locally of finite type, hence $s', t'$ locally of finite
presentation (because $U'$ is a disjoint union of spectra of fields
in particular locally Noetherian, so that
Morphisms of Spaces, Lemma
\ref{spaces-morphisms-lemma-noetherian-finite-type-finite-presentation}
applies). Hence $\mathcal{Z} = [U'/R']$ is an algebraic stack by
Criteria for Representability, Theorem
\ref{criteria-theorem-flat-groupoid-gives-algebraic-stack}.
As $R'$ is the restriction of $R$ by $U' \to U$ we see
$\mathcal{Z} \to \mathcal{X}$ is a monomorphism by
Groupoids in Spaces, Lemma
\ref{spaces-groupoids-lemma-quotient-stack-restrict}
and
Properties of Stacks, Lemma \ref{stacks-properties-lemma-monomorphism}.
Since $\mathcal{Z} \to \mathcal{X}$ is a monomorphism we see that
$|\mathcal{Z}| \to |\mathcal{X}|$ is injective, see
Properties of Stacks, Lemma
\ref{stacks-properties-lemma-monomorphism-injective-points}.
By
Properties of Stacks, Lemma \ref{stacks-properties-lemma-points-cartesian}
we see that
$$
|U'| = |\mathcal{Z} \times_\mathcal{X} U'|
\longrightarrow
|\mathcal{Z}| \times_{|\mathcal{X}|} |U'|
$$
is surjective which implies (by our choice of $U'$) that
$|\mathcal{Z}| \to |\mathcal{X}|$ has image $\{x\}$.
We conclude that $|\mathcal{Z}|$ is a singleton.
Finally, by construction $U'$ is locally Noetherian and reduced, i.e.,
$\mathcal{Z}$ is reduced and locally Noetherian. This means that
the essential image of $\mathcal{Z} \to \mathcal{X}$
is the residual gerbe of $\mathcal{X}$ at $x$, see
Properties of Stacks, Lemma
\ref{stacks-properties-lemma-residual-gerbe-unique}.
\end{proof}

\end{document}
